\chapter{The Veil and the Cradle}

\section*{Told by}
The Vizier's widow, a woman who wears a mask even inside her own home. She speaks through a silk curtain. No one has seen her face in twenty years---not her servants, not her physicians, not the priests who came to bless her late husband's tomb. The mask has seven latches. She has not opened them since the night the palace fell. She says the face beneath is not the one she was born with. She says it changed. She does not say how.

\section*{The Tale}

A slave girl was bought for the magistrate's harem. She was young---too young, the other women whispered---frightened, and silent. Not the silence of patience. The silence of a throat that has forgotten how to make sound. She had been taken from a village that no longer existed, from a family whose names she had been beaten into forgetting. The magistrate's wife took pity on her. Pity in a harem is a dangerous thing. It is also the only honest thing.

\textit{``Wear a veil,''} the wife said. \textit{``The magistrate will not see you. He sees only the silk, only the pattern, only the suggestion of a face. You may become invisible. You may become safe.''}

The girl wore the veil. It was pale blue, the color of a winter sky, embroidered with silver thread in patterns that shifted when she moved. She served the magistrate's table. He did not see her. She cleaned his chambers. He did not remember her. She walked the halls of the palace like a ghost---present, but never noticed, never questioned, never quite there.

For three years, she lived like this. She learned the palace's secrets: which guards slept on duty, which passages led to the garden, which courtiers were loyal and which were waiting. She could have killed the magistrate a dozen times. She did not. She was not a knife. She was a witness.

One night, the magistrate's enemy sent a knife. Not a metaphor---a blade, narrow as a finger, poisoned with something that turned the edge black. The assassin came through the window, silent as smoke, and raised the blade over the magistrate's sleeping throat.

The girl saw it. She could not call out---her voice would break the veil, and a broken veil is just a piece of cloth. She stepped between the knife and the magistrate. Not thinking. Not choosing. Simply moving, the way water moves toward a crack in a dam.

She took the blade in her chest. The veil fell. The magistrate woke to the sound of silk hitting stone.

He saw her for the first time. Her face was young---younger than he had expected---and her eyes were the color of the veil, winter sky, something distant and cold.

\textit{``Who are you?''} he asked.

She opened her mouth. After three years of silence, she had almost forgotten how to shape words. But she found them.

\textit{``The one who was never there,''} she said.

Then she died. The poison was fast. The magistrate did not have time to call for a healer. He did not have time to learn her name. He did not even have time to close her eyes.

The magistrate had the veil woven into a shroud for her---blue silk and silver thread, the same veil that had hidden her, now the only thing that would touch her skin. He kept the shroud for the rest of his life, folded in a cedar chest beside his bed. He never opened the chest. He never married again. He told his advisors that he had seen the face of mercy, and mercy had not worn silk. Mercy had worn a slave's collar and a blade in her chest.

The assassin was never found. But the magistrate's enemies learned that a veil is not just a hiding place. Sometimes it is a cradle. And the ones who are invisible are the ones who see the most clearly.

\section*{The Witch's Closing}
\textit{``The veil hides. The knife reveals. The cradle holds what cannot be seen. Choose which you carry---but know that the one who carries the veil may also carry the weight of every face she has hidden from. The silence is not empty. The silence is full of names.''}

\section*{What Lurks in the Shadow of the Tale}

\subsection*{The Silent Veil}

A silk veil, pale blue with silver thread, that makes the wearer unnoticeable---not invisible, but forgettable. Guards will look past her. Servants will not remember her face. Courtiers will speak freely in her presence, not because they trust her, but because they do not see her at all. She is a gap in perception, a blind spot in the architecture of attention.

She cannot speak while wearing it, or the effect breaks. Not gradually. Not with a warning. The moment her voice leaves her throat, every eye in the room will turn to her. Every guard will remember the face they did not see. Every door will close.

\paragraph{Tags / Mechanics:}
\begin{itemize}
\item [STEALTH] [CURSE] [SILENCE]
\item While wearing the veil, the character gains +2 dice to Stealth and Deception to avoid notice. She may walk through crowds, past guards, across courtyards, as long as she does nothing to draw attention. She may open doors, pick locks, take objects---as long as she is quick and quiet. But the moment she speaks, the veil becomes just a piece of cloth.
\item The veil does not hide sound. Footsteps, breathing, the rustle of fabric---these can still be heard. A guard listening at a door will know someone is there. He just will not remember who.
\item Each scene, the GM may ask the character to test Spirit + Resolve (DV 4). On failure, the character feels compelled to speak---a gasp, a cough, a whispered prayer. The compulsion is not magical. It is the weight of three years of silence pressing against the inside of her ribs. On failure, the veil's effect ends for the scene, and nearby creatures are alerted.
\item Removing the veil is an action. Replacing it is also an action. The veil does not forgive interruption. Once broken, it takes a full day of wearing it in silence for the effect to return.
\end{itemize}

\paragraph{Adventure Hook:}
A noble's daughter has been sold into a hostile court---a marriage she did not choose, to a man she has never met, in a city where she has no allies. She possesses a veil that makes her unnoticeable, but she cannot speak while wearing it. The party must smuggle her out of the court without being seen, without raising the alarm, and without her breaking the silence. But the veil has a cost: every hour she wears it, she forgets another face. First her mother's. Then her father's. Then her own. By the time she reaches the gate, she may not remember why she wanted to leave.

\section*{Colophon}
A room behind a curtain. No window. No second door. The widow poured tea into a cup she could not see---her hands knew where the cup was, had known for twenty years. The cup rattled against the saucer. \textit{``I wore it for twenty years,''} she said. \textit{``Now I do not need it. The silence is inside me. It has always been inside me. The veil was just a way of keeping it company.''} She did not drink the tea. The tea grew cold. The curtain did not move.